%! Bias in Samples and Surveys — Unit 1, Lesson 1.4 — Lesson Plan
% PILOT SHAPE (2026-09-06): modeled on AP Statistics lesson 1.4 minus its AP Practice page.
% GRADUAL-RELEASE plan (warm-up / guided notes and practice / spoken debrief / close and START
% THE HOMEWORK). The guided notes carry the release: I do (two sections) -> we do (Guided
% Practice). THERE IS NO INDEPENDENT PRACTICE SET IN THE NOTES — the homework is the individual
% practice, and the period ends with students starting it. Teacher-facing. Breakable boxes are
% `skillbox` with a \boxguard ahead; the two tabularx boxes are `fixedskillbox` (never splits).
% No group activity, no tier boxes, no exit ticket.
\documentclass[10pt]{article}
\usepackage{saar-article}
\usepackage{saar-boxes}
\usepackage{graphicx}   % -article does NOT load graphicx
\graphicspath{{images/}}
\newcommand{\TallMath}[1]{$\displaystyle #1\rule[-1.4em]{0pt}{3.2em}$}
\newcommand{\UnitNumberName}{Unit 1: Asking Questions with Data: Study Design \quad}
\newcommand{\LessonNumberName}{Lesson 1.4: Bias in Samples and Surveys}

\begin{document}
\begin{center}
    {\Large\bfseries \CourseName \\
    {\small\bfseries \UnitNumberName \LessonNumberName}}
\end{center}

\begin{tcolorbox}[colback=frost, colframe=cerulean, boxrule=0.9pt, arc=2mm,
    left=2mm, right=2mm, top=1.5mm, bottom=1.5mm]
    \textbf{Primary Objective:} Students will explain what makes a study \emph{biased} rather
    than merely variable; identify sampling bias in a plan, naming \emph{undercoverage} and
    \emph{nonresponse} by asking whether the missing person could have been chosen; name the
    forms of response bias and say which direction each pushes an estimate; show by
    computation that a larger sample does not reduce bias; and describe and justify a change
    that removes a named bias from a plan.
    \\[2pt]
    \textbf{Standards:} PS.DC.2c; AFDA.DA.2e, AFDA.DA.2f, AFDA.DA.2g \quad$\cdot$\quad
    PS.DC.1d--e and PS.DC.2a carried forward in the population/sample/who-chose items
    \quad$\cdot$\quad previews PS.DC.2d (Lesson 1.5)
    \\[2pt]
    \textbf{Lesson model:} \emph{Gradual release --- I do, we do, you do --- all of it inside
    the guided notes.} There is no group activity. The notes name the vocabulary as they teach
    it and deliver the instruction in two sections; the Guided Practice is worked together with
    students holding the pen; \textbf{the homework is the individual practice}, and students
    start it in class while you circulate. The whole-class debrief is \emph{spoken} --- there is
    no practice set at the end of the notes and no exit ticket.
    \\[2pt]
    \textbf{Note on the ten-week study:} the crux depends on students believing that $750$
    members beats $75$. Do not compute, hint at, or resolve the ten-week result until the
    second half of section~2 --- the hook and section~1 run on the $75$ alone.
\end{tcolorbox}

\renewcommand{\arraystretch}{1.4}
\boxguard
\begin{skillbox}[Learning Targets \& Key Understandings]{goldbox}
\begin{tabularx}{\linewidth}{@{}X|X@{}}
\textbf{Learning targets (student language)} & \textbf{Key understandings (the ``why'')} \\[2pt]
\hline
\vspace{-6pt}
\begin{itemize}[leftmargin=1.1em, nosep, after=\vspace{2pt}]
  \item I can explain what makes a study \textbf{biased} and tell bias apart from
        \textbf{sampling variability} --- a direction, not bad luck. \hfill \textit{PS.DC.2c}
  \item I can identify \textbf{sampling bias} in a plan, naming \textbf{undercoverage} and
        \textbf{nonresponse}, by asking \emph{could that person have been chosen?}
        \hfill \textit{PS.DC.2c; AFDA.DA.2f}
  \item I can name the forms of \textbf{response bias}, including \textbf{question order
        bias}, and say which \textbf{direction} a plan pushes the estimate.
        \hfill \textit{PS.DC.2c; AFDA.DA.2e}
  \item I can explain why a \textbf{larger sample does not fix bias}, and describe and
        \emph{justify} a change that removes it. \hfill \textit{PS.DC.2c; AFDA.DA.2g}
\end{itemize}
&
\vspace{-6pt}
\begin{itemize}[leftmargin=1.1em, nosep, after=\vspace{2pt}]
  \item \emph{Bias is a direction, not bad luck.} Lesson 1.2's honest samples scattered
        \emph{around} the parameter; a biased procedure misses the same way every time. Bias is
        a property of the \textbf{method}, not of any one sample.
  \item \textbf{Target misconception:} that ``the sample was too small'' --- that more of the
        same data drifts toward the truth on its own. It does not: a larger sample makes an
        estimate more \emph{precise}, not more \emph{correct}. Size shrinks variability and
        does nothing to direction; ten Tuesdays give $36\%$ exactly as one did. The standard
        says it outright: \emph{a large sample size does not make up for bias.}
  \item \textbf{Second misconception:} undercoverage and nonresponse fused as ``people were
        left out.'' The split is whether they were ever \emph{on the frame} --- \emph{could that
        person have been chosen?} No $\Rightarrow$ undercoverage. Yes, but silent
        $\Rightarrow$ nonresponse. The $200$ returned cards have a perfect frame and a broken
        study.
  \item Lesson 1.3 chose \emph{who} to ask. Response bias happens \emph{after} the right
        person has been chosen --- who asks, how it is worded, what order it comes in --- so no
        sampling method prevents it.
  \item Naming a bias is half of AFDA.DA.2g; \textbf{attaching the fix to what it removes} is
        the other half. ``Ask better questions'' is not an answer; ``a complete frame removes
        undercoverage'' is.
\end{itemize}
\end{tabularx}
\end{skillbox}

\boxguard[24]
\begin{skillbox}[{Vocabulary, Concepts, \& Theorems --- taught directly in the guided notes}]{greenbox}
\small
\renewcommand{\arraystretch}{1.35}
\begin{tabularx}{\linewidth}{@{} >{\bfseries}p{3.1cm} X @{}}
Bias & A study is biased when its procedure systematically overestimates or underestimates
       the population value --- every sample it takes leans the same direction. \\
Sampling variability & Honest samples disagree with each other and scatter \emph{around} the
       truth. Noise, not error; it shrinks as the sample grows. \\
Sampling bias & Bias from how the sample was chosen: some members of the population are
       systematically more likely to be picked than others. \\
Undercoverage & Part of the population is left off the sampling frame, so those members
       have no chance at all of being chosen. \\
Nonresponse & Individuals are chosen properly but never answer, so the data is missing
       exactly the people who stayed silent. \\
Response bias & Anything about the survey itself --- who asks, the wording, the order --- that
       pushes a respondent away from their true answer. \\
Forms of response bias & Demand, social desirability, dissent, acquiescence, extreme
       responses, neutral responding, question order (the standard's seven). \\
Question order bias & An earlier question changes the answer to a later one, so reordering the
       same questions changes the result. \\
Precise vs.\ correct & A larger sample buys precision (less scatter). Only a better procedure
       buys correctness (no lean). \\
\end{tabularx}
\end{skillbox}

% fixedskillbox never splits, so the phase table stays intact on one page.
% THE PHASE TABLE IS A CONTRACT: 5 / 35 / 10 / 10 — the I do is ~20 of the 35 and the Guided
% Practice ~15; the debrief is spoken; the homework start is the second formative read.
\begin{fixedskillbox}[Lesson at a Glance --- \MeetingLength]{frost}
\renewcommand{\arraystretch}{1.25}
\begin{tabularx}{\linewidth}{@{}l c X X@{}}
\textbf{Phase} & \textbf{Min} & \textbf{Students} & \textbf{Teacher} \\
\hline
Warm-Up & 5 & Harbor Point's $1{,}500$: the census percent, the Tuesday sample's percent, the
  prediction against the truth --- silent & Debrief item~3 aloud; write ``$20\%$ \; $36\%$ \;
  $540$ vs.\ $300$ --- off by $240$, too high'' on the board and leave it up --- section~1
  opens on it \\
Guided Notes \& Practice & 35 & Fill the vocabulary box and notes 1--2 as each term is named;
  then \emph{Cedar Ridge Apartments} together, pen in their hand & \textbf{I do} notes 1--2
  (20) $\to$ \textbf{we do} the Guided Practice box (15) \\
Debrief & 10 & Pens down, papers up --- answer aloud, nothing to fill in & The ten Tuesdays
  back on the board; the ``too small'' reflex, then the ``left out means undercoverage''
  reflex; the Bayside cold check \\
Close \& Assign & 10 & \textbf{Start the homework in class}, alone and silent & Announce the
  homework and its form (packet), circulate for the second formative read, preview
  Lesson~1.5 \\
\end{tabularx}
\end{fixedskillbox}

\begin{fixedskillbox}[Warm-Up --- Activate Prior Knowledge (5 min)]{frost}
\begin{tabularx}{\linewidth}{@{}X X@{}}
\begin{minipage}[t]{\linewidth}\vspace{0pt}
    \textbf{The three items and what each seeds}
    \begin{itemize}[leftmargin=1.1em, itemsep=1pt, topsep=2pt]
      \item \textbf{1.} $300$ of $1{,}500$ as a percent. \textbf{Spirals} to Lesson 1.1's
            relative frequency; \textbf{seeds} the \emph{truth} column --- $20\%$ is the
            parameter every table in the notes is measured against.
      \item \textbf{2.} $27$ of $75$ as a percent. \textbf{Seeds:} the Tuesday-evening
            statistic, $36\%$ --- the study 1.3 threw out, now the biased row of section~1's
            table and the ten-week study's unchanging answer.
      \item \textbf{3.} $0.36 \times 1500 = 540$, off by $240$, too high. \textbf{Seeds:} the
            first time in Unit~1 a statistic sits beside a \emph{known} parameter --- the
            \emph{direction} of the miss is the whole of section~1's first half.
    \end{itemize}
\end{minipage}
&
\begin{minipage}[t]{\linewidth}\vspace{0pt}
    \textbf{Running it}
    \begin{itemize}[leftmargin=1.1em, itemsep=1pt, topsep=2pt]
      \item \textbf{Debrief item~3 aloud} and write $20\%$, $36\%$, $540$ vs.\ $300$, \emph{off
            by $240$, too high} on the board. Leave it up all period; section~1 hangs the word
            \emph{bias} on that line, and the hook is asked against it.
      \item Item~3(b) is the tell. Students who write $16$ instead of $240$ answered in
            percentage points, not members --- both are right ideas; the points form is exactly
            what section~2's table wants, so credit it and tell them to hold the thought.
      \item Say that all three are Algebra~1 arithmetic and are not being retested. \textbf{Do
            not say the word \emph{bias} here, and do not explain why the prediction misses}
            --- the notes name what the class has already computed.
    \end{itemize}
\end{minipage}
\end{tabularx}
\end{fixedskillbox}

\boxguard
\begin{skillbox}[Hook --- before anyone sits down]{frost}
The hook is on \textbf{page 1 of the notes}, under the vocabulary box, and on the board:
\textit{``The Tuesday-evening study said $36\%$; the records say $20\%$. The staff member says
the problem was size --- she will run the same Tuesday-evening survey for ten weeks and ask
$750$ members instead of $75$.''} \textbf{Ask: will her new number land closer to the truth?}
Students circle \emph{yes}, \emph{no}, or \emph{it depends} and write one line of reason in the
hook box \emph{before} anyone speaks; then take the hand vote and write the split on the board.
Most rooms vote \emph{yes} --- ``more people, better answer'' is the reflex a year of Algebra~1
has taught them. \textbf{Do not settle it, and do not compute anything about the $750$.} Say
only that the answer is the last thing in section~2 and that they will vote again then. The
vote is what makes the crux land: students have to have committed to \emph{bigger is closer}
before the notes take it away.
\end{skillbox}

% The multicols body is one unsplittable block about a page long; a small guard lets the two
% opening lines print as a stub at the foot of the page above (seen 2026-09-07). [40] forces it over.
\boxguard[40]
\begin{skillbox}[Guided Notes \& Practice --- I do, then we do (35 min)]{frost}
\textbf{This is the whole lesson.} There is no group activity, and there is no independent
practice set on this page: \textbf{the homework is the individual practice}, started in class.
\begin{multicols}{2}
\textbf{I do (about 20 min).} Two sections, each in two moves; students fill the blanks as each
idea is named, and the vocabulary box on page~1 is filled \emph{as you go}, never in advance.
\textbf{Every fill-in cluster opens with a word bank}; the work is choosing and justifying,
not recalling.

\emph{Section 1} opens on the two-row table, which does the conceptual job visually. Both rows
wobble; only one is wrong the same way every time. \textbf{Ask before revealing the averages:
which row would you trust?} Students reliably pick the tighter Tuesday row --- ``it agrees with
itself'' --- and that wrong answer is worth two minutes: \emph{consistent is not correct}. Fill
the averages, name \textbf{bias}, and land the direction: too high, because Tuesday evening is
lap-swim time. Five minutes. Its second half, \emph{Sampling Bias --- Who Never Had a Chance},
is 1.3's frame idea wearing a new name. Run the four-row table fast and \textbf{save the
argument for row~4}: all $1{,}500$ got a card, the frame is perfect, and the study is still
broken. Ask what is different about the $200$ who mailed theirs back. That is nonresponse, and
the red callout's test --- \emph{could that person have been chosen?} --- is the line to make
them say aloud. Five minutes.

\textbf{Section 2 is the lesson.} Its first half is response bias: read the seven-row reference
table \emph{once} --- do not lecture it --- then go straight to the four diagnostic rows and
let the class argue. Rows 2 and~3 (neutral responding vs.\ acquiescence) are the pair to slow
on. Then the question-order sentence: identical words, $21$ points apart. Six minutes. Its
second half, \emph{The Fix That Is Not a Fix}, is the crux. Work $270 \div 750 = 36\%$ with
the room, fill the comparison table --- $16$ points, then $16$ points again --- and
\textbf{re-take the hook vote before saying a word}. Then land the blue callout, blank by
blank: \emph{precise, not correct; variability, not direction}; and the three fixes, each
attached to the bias it removes. Four minutes.

\columnbreak
\textbf{We do (about 15 min).} The Guided Practice box, \emph{Cedar Ridge Apartments}, alone on
page~4 --- the same moves the homework will ask alone, on an apartment complex instead of a
pool. Read the setup aloud when you open it: \emph{every} household got a card --- part (b)
does not work if students skim past that word. \textbf{Students hold the pen; you hold the
questions.} Part (a) is warm-up arithmetic with two different denominators on purpose ---
$150/600$ and $96/150$; a student who divides $96$ by $600$ has merged the two questions.
Part (b) is the red callout transferred. Part (c) is the scale-and-compare move from the
warm-up: $384$ against $240$, $24$ points, too high. \textbf{Part (d) is the crux
transferred}: twice the cards, $64\%$ again, and the car owners answered at twice the rate.
Take a hand count on \emph{did the extra cards help?} before anyone writes, and make both sides
defend it. Part (e) is AFDA.DA.2g --- the wording names the bias, and the rewrite has to be
genuinely balanced, not leaning the other way.

Circulate during Guided Practice and demand the reason, not the word:
\begin{itemize}[leftmargin=1.1em, nosep]
  \item \textbf{Part (a):} ``Ninety-six out of how many? Who is in that denominator?''
  \item \textbf{Part (b):} ``Could a silent household have been chosen? Then which bias?''
  \item \textbf{Part (c):} ``Too high or too low --- and who is over-counted?''
  \item \textbf{Part (d):} ``If they mailed to every household ten times, what changes?'' A
        student who says \emph{they'd get closer} has not left the hook vote.
  \item \textbf{Part (e):} ``Read me your rewrite. Which way does \emph{it} lean?''
\end{itemize}
While circulating, pick the papers to put up in the debrief --- one part (d) that said
\emph{same people, more of them} and one that said \emph{still not enough cards}. \textbf{Your
formative read comes twice}: here, and again in the last ten minutes as students start the
homework.
\end{multicols}
\end{skillbox}

\boxguard[24]
\begin{skillbox}[Debrief --- whole class, spoken (10 min)]{frost}
Pens down, papers up. \textbf{This debrief is spoken} --- there is no box at the end of the
notes to fill in, so it lives entirely on the board and in the room.
\begin{multicols}{2}
\begin{enumerate}[leftmargin=1.3em, nosep, itemsep=3pt]
  \item \textbf{Put the ten Tuesdays back up} and make a student say the sentence: $750$ is
        not closer than $75$, because size shrinks variability and this study's problem is
        direction. Then state it as the standard does: \emph{a large sample size does not make
        up for bias.}
  \item Put up the two Guided Practice part (d)s you picked. The difference between them is
        not arithmetic --- it is whether the student still believes more data drifts toward
        the truth on its own.
  \item Circle the $200$ returned cards and take the \emph{left-out-means-undercoverage}
        reflex: \emph{``Could a silent card holder have been chosen? Then which bias is it?''}
  \item Take part (e) last, and insist on balance: a rewrite that leans the other way
        (``Is there plenty of parking?'') has learned that wording matters but not that
        \emph{balance} is the fix.
\end{enumerate}
\columnbreak
\textbf{Check for understanding --- ask it cold, whole class.} \emph{``Bayside Middle School,
$750$ students. The principal asks $30$ students in the library at lunch whether lunch is too
short: $24$ say yes --- $80\%$. Told her sample was too small, she asks $60$ library students:
$48$ say yes --- $80\%$ again. Name the bias in her plan. Did the larger sample fix it? Why
not?''} Hands down, thirty seconds of think time, then cold-call three students. Correct
answer: \textbf{undercoverage} --- only students in the library were on the frame; a student in
the cafeteria or outside could never have been chosen --- and \textbf{no}: the $60$ are still
library students, so the lean is the same; size touches variability, not direction.
\emph{This replaces the exit ticket.}

\textbf{Formative read.} Sort the three answers as they come: \textbf{(a)} named
undercoverage and said size does not touch direction; \textbf{(b)} named it but said $60$ is
``still too small''; \textbf{(c)} said the larger sample fixed it, or called it nonresponse.
Pile (b) is the teachable one --- it is the right instinct with the wrong lever (the procedure,
not the count). \textbf{If pile (c) runs past a third of the room, open Lesson~1.5 by
re-running section~2's second half}, and say so plainly.

\textbf{If the debrief runs short}, cut items 3 and~4 --- item~1 is the one that cannot be
cut. \textbf{Do not let the debrief eat the last ten minutes} --- the homework start is not
optional padding, it is your second formative read.
\end{multicols}
\end{skillbox}

\boxguard
\begin{skillbox}[Homework --- scored, started in class, due after two study halls]{goldbox}
Six exercises on \textbf{Millbrook Public Library} ($1{,}200$ card holders; a front-desk pile
that says $80\%$ and a proper phone sample that says $45\%$): \textbf{1} population, sample,
\emph{convenience}, undercoverage, and who had no chance (\emph{the spiral to Lessons 1.2 and
1.3}); \textbf{2} $72/90 = 80\%$ scaled to $960$ --- the core procedure; \textbf{3} the phone
sample, $54/120 = 45\%$ scaled to $540$, and the gap in points and in card holders ---
interpret in context; \textbf{4} Plan A (four weeks at the desk, $360$ back) against Plan B
(the phone sample) --- \textbf{the contrast pair}, bigger vs.\ closer --- then what the $360$
would have reported: \emph{the crux head-on}; \textbf{5} four response biases to name and the
$21$-point order effect; \textbf{6} a random draw with a $65\%$ response rate --- explain why
it can still be biased and name the fix (AFDA.DA.2g). Capped at \textbf{two pages} (user
decision 2026-09-06).

\textbf{Start it in the last ten minutes of the period, in class and alone.} \textbf{Scoring:}
12 points --- 2 per item. \textbf{Item~4 is where the misconception shows}: a student whose
$360$ report anything near $45\%$ believes more data drifts toward truth on its own.
\textbf{Item~6 carries the second check}: ``$78$ is too small'' is the wrong lever again; the
answer has to be that a random draw decides who is \emph{asked}, not who \emph{answers}.

\textbf{Packet or Desmos --- decided here.} The packet pages are authored and are the default.
Desmos Classroom has sampling simulations, but nothing that asks a student to diagnose a
plan, name the bias, say its direction, and attach a fix --- so \textbf{1.4 is a packet
night}. Say so aloud at Close \& Assign.

\textbf{Preview of Lesson~1.5.} Every study in Unit~1 so far has only \emph{watched}. Next
class names that: observational studies, what they can and cannot support (PS.DC.2d). Do not
open causation here.

\textbf{Connections \& Big Ideas.} \textit{Unit 1 core idea:} a conclusion is only as
trustworthy as the study that produced it. \textit{Standards covered:} PS.DC.2c; AFDA.DA.2e,
AFDA.DA.2f, AFDA.DA.2g. \textit{Spirals forward to:} observational studies (1.5, PS.DC.2d), experimental design (1.6, PS.DC.3a--b), the
full-cycle project (1.8), and --- directly --- the conditions on inference in Units 7 and~8, where
a biased sample invalidates an interval no matter how narrow. \textit{Reaches back to:} Lesson
1.2's sampling variability and Lesson 1.3's sampling frame.
\end{skillbox}

\boxguard
\begin{skillbox}[Watch For (while circulating)]{redbox}
\begin{itemize}[leftmargin=1.2em, nosep]
  \item \textbf{``The sample was too small''} (notes~2 second half, GP~(d), HW~4 --- the target
        misconception). This is the reflex answer to every flawed study in the lesson and it is
        wrong in all of them. Probe: \emph{``If they asked ten times as many the same way, what
        would change?''}
  \item \textbf{Undercoverage and nonresponse fused} (notes~1 row~4, GP~(b), HW~1 and~6). Both
        leave people out; the split is whether they were ever on the frame. Probe: \emph{``Could
        that person have been chosen?''} Yes but silent $\Rightarrow$ nonresponse. No
        $\Rightarrow$ undercoverage.
  \item \textbf{``Biased'' used to mean ``prejudiced''} (everywhere). Students arrive with the
        everyday meaning and write ``the survey was unfair.'' Redirect every time:
        \emph{``Which direction does the number miss, and by how much?''}
  \item \textbf{A bias named with no direction} (notes~1, GP~(c), HW~3). Probe: \emph{``Too high
        or too low --- and who is over-counted?''} The label alone is half credit.
  \item \textbf{Neutral responding read as extreme, or as acquiescence} (notes~2 diagnostic
        rows, HW~5). Probe: \emph{``Which boxes did that person actually circle?''}
  \item \textbf{A fix with no bias attached} (notes~2 callout, GP~(e), HW~6). ``Ask better
        questions'' and ``do it properly'' are not answers. Probe: \emph{``What does your fix
        remove?''}
\end{itemize}
Cold-call prompts: \emph{``Could that person have been chosen?''} \quad \emph{``Which way does
it push?''} \quad \emph{``What would ten times as many change?''} \quad \emph{``What does your
fix remove?''}
\end{skillbox}

\boxguard
\begin{skillbox}[Close \& Assign (10 min)]{goldbox}
\textbf{Students start the homework here, in class, alone and silent --- this is the individual
practice, and the first minutes of it are your best formative read.} Hand the launch first ---
six exercises on paper, scored out of 12, due at the start of the first class after two study
halls --- and point out that the \emph{Keep in Mind} box on the cover is the page to flip back
to. Circulate: \textbf{item~4 is the column that tells you whether section~2 landed}, and it is
on page~2, so put students on 1, 4, and~6 first and let 2, 3, and~5 go home. Sort what you see
into three piles as you walk --- wrote \emph{about $80\%$} with \emph{same desk, same people}
(ready); wrote \emph{about $80\%$} with no reason (ready, needs the language); wrote something
near $45\%$, or ``$360$ is enough'' (the misconception survived) --- and open Lesson~1.5 by
re-running section~2's second half if that third pile ran past a third of the room. Then close
on what changed: \emph{``You walked in able to scale $36\%$ to $540$. You walk out knowing that
a study can be wrong the same way every time, that no sample size fixes that, and that the only
repair is a better procedure --- a complete frame, a call back to the silent, a question that
does not lean.''} \textbf{Preview:} Lesson~1.5, \emph{Observational Studies} --- every study so
far has only watched; next time we ask what a study that only watches can and cannot prove.
\end{skillbox}

% ── Teacher notes — one per component, in packet order (three) ──────────────
% Teacher-only prose lives HERE, never in a _key: a note is the one block in a key with no
% counterpart in the blank. No note for the debrief or the homework start — both are fully
% specified in their own boxes above.
\begin{teachernote}[Warm-Up]
Five minutes, silent, timed. Items 1 and~2 are Algebra~1 percent work and should be near
universal --- say so when you assign it. \textbf{Item~3 is the one that matters}: $0.36
\times 1500 = 540$, against the $300$ the records say --- off by $240$ members, too high.
This is the first time in Unit~1 that a sample statistic is put beside a \emph{known}
parameter. Do not explain why yet, and do not say the word \emph{bias}. Write $20\%$, $36\%$,
$540$ vs.\ $300$, \emph{off by $240$, too high} on the board and \textbf{leave it up all
period}; section~1 points back at it, and the hook is asked against it. Watch for students who
answer 3(b) with $16$ instead of $240$: they answered in percentage points, not members. Both
are right ideas, and the points form is exactly what section~2's comparison table will want ---
credit it and tell them to hold the thought.
\end{teachernote}

\begin{teachernote}[Guided Notes \& Practice]
\textbf{This one document is the lesson --- pace it 20 / 15, then 10 for the debrief, then 10
for the homework start.} Section~1 must move: five minutes on the two-row table and the word
\emph{bias}, five on undercoverage and nonresponse. \textbf{Ask which row they would trust
before revealing the averages}; if the averages go up first, the trap is gone and
\emph{consistent is not correct} becomes a sentence to copy. Row~4 of the second table is the
one that argues --- every member got a card, so the frame is perfect, and the study is still
wrong. If students insist that is the same as undercoverage, use the test question ---
\emph{could that person have been chosen?} --- and hold the line; the Guided Practice's entire
premise is this row.

\textbf{Spend the time in section~2}, which carries the misconception. The seven response
biases are recognition-level, not recall-level: read the reference table once and go to the
four diagnostic rows. The pair to slow down on is neutral responding versus acquiescence,
which get fused into ``didn't take it seriously.'' Then the crux. Work $270 \div 750$ with the
room, fill the comparison table --- $16$ and $16$ --- and \textbf{re-take the hook vote before
you say a word}. Let the room argue; the students who voted \emph{yes} will reach for
\emph{maybe ten weeks still isn't enough}. Ask them what a hundred weeks would change. Then
land the blue callout blank by blank and have the class say the sentence aloud before they
write it: \emph{a larger sample makes an estimate more precise, not more correct.} This is the
single most transferable line in Unit~1 and it returns in Unit~7 as the difference between a
narrow interval and a right one.

\textbf{Guided Practice is where the pen changes hands.} \emph{Cedar Ridge Apartments} runs the
moves the homework will ask alone --- two rates with different denominators, the frame test,
scale and compare with a direction, the crux, a fix attached to a bias --- on an apartment
complex instead of a pool. Read the setup aloud and lean on the word \emph{every}. Part~(a)'s
two fractions have different denominators on purpose; a group that gets $16\%$ divided $96$ by
$600$. Part~(d) is the crux transferred: twice the cards, $64\%$ again, and the direction
follows from the car owners answering at twice the rate. A student who defends it with ``they
still need more cards'' has not left the hook vote. Part~(e)'s rewrite must be balanced, not
flipped --- ``Is there plenty of parking?'' leans the other way. If you only get through (a)
through (d), that is the right trade; (e) can be said aloud.

\textbf{There is no independent practice set on this page, and that is deliberate --- the
homework is the individual practice.} Guided Practice is the last thing on the page; the period
ends with students starting the homework in front of you, which is where your second formative
read comes from. Item~4 is the one to watch. If the ``$360$ is enough'' pile is more than a
third of the room, \textbf{open Lesson~1.5 by re-running the second half of section~2}, and say
so plainly.

\textbf{Debrief (10 min), spoken.} The ten Tuesdays back on the board with a student saying
\emph{$750$ is not closer, because the problem is direction}; it is the one move that cannot
be cut. \textbf{Do not let the debrief eat the last ten minutes} --- the homework start is not
optional padding.
\end{teachernote}

\begin{teachernote}[Homework]
\textbf{Scored, 12 points (2 per item), due at the start of the first class after two study
halls.} The ten minutes at the end of the period are a \emph{start}, not a completion: put
students on 1, 4, and~6 while you can still catch a wrong start, and let 2, 3, and~5 go home.
Items 2 and~3 are the spine and are designed to be done in order: $80\%$ and $960$, then $45\%$
and $540$, then the gap --- $35$ points and $420$ card holders. A student who computes the gap
both ways has understood that percentage points and counts are two views of the same error.
\textbf{Item~4 is the one to read first when scoring} --- the answer is a number (\emph{about
$80\%$ again}) plus a reason (\emph{same desk, same visitors, more of them}). A student who
predicts anything near $45\%$ believes more data drifts toward truth on its own; that is the
misconception to correct in writing, not just mark. Item~5's rows 1 and~4 are the pair students
fuse --- agreeing with everything is acquiescence; the middle box every time is neutral
responding. Item~6 closes the arc the hook opened: the draw was perfectly random and the counts
are all correct, so the answer cannot be ``they sampled badly'' or ``$78$ is too small.'' It has
to be that a random draw decides who is \emph{asked}, not who \emph{answers} --- the $42$ who
never picked up may use the library differently --- and the fix is to keep calling back.
Score items 1--3 and~5 for correctness and items 4 and~6 for the \emph{reason}, not the verdict.
\textbf{When you score the page, sort item~4 into three piles} --- said \emph{same procedure,
same lean} (ready); said \emph{about $80\%$} with no reason (ready, needs the language); said
something near $45\%$ (the misconception survived). If more than a third land in the last
pile, reopen the ten Tuesdays in the Lesson~1.5 warm-up debrief rather than letting it ride.

\textbf{Packet is the call for 1.4.} Desmos carries sampling simulations and none of the
diagnose-name-direction-fix reasoning in this set; do not swap it in. The \textbf{spiralbox}
hands off to Lesson~1.5 deliberately: every study in this unit so far has only watched. Do not
open causation here.
\end{teachernote}
\end{document}
